\documentclass{article}
\usepackage[english]{babel}
\usepackage{geometry,amsmath,amssymb,amsthm,hyperref}
\geometry{letterpaper, margin=1.2in}

\newtheorem{theorem}{Theorem}
\newtheorem{lemma}{Lemma}
\newtheorem{proposition}{Proposition}
\newtheorem{definition}{Definition}
\newtheorem{corollary}{Corollary}
\newtheorem{remark}{Remark}

\newcommand{\C}{\mathbb{C}}
\newcommand{\R}{\mathbb{R}}
\newcommand{\N}{\mathbb{N}}
\newcommand{\Psimu}{\Psi_{\mu}}
\newcommand{\Imu}{I^{\mu}}
\newcommand{\Dmu}{D^{\mu}}

\hypersetup{
  colorlinks=true,
  linkcolor=blue,
  citecolor=blue,
  urlcolor=blue
}

\begin{document}

\title{%
  Fox--Wright Orthogonal Polynomials and the\\
  Fractional Riccati--M\"{u}ntz--Pad\'{e} Theorem:\\
  A New Family in Non-Hermitian Approximation Theory
}
\author{Stephen Crowley}
\date{May 2026}
\maketitle

\begin{abstract}
We identify and characterize the orthogonal polynomial sequence canonically
associated with the Fractional Riccati--M\"{u}ntz--Pad\'{e} theorem
\cite{Crowley2026}. The polynomials arise as the denominators of the
Pad\'{e} approximants built from the Müntz--Tau moment sequence
$\{a(k,u)\}_{k \ge 0}$ of the fractional Riccati equation
$D^{\mu}y = P(u) + Q(u)y + R(u)y^{2}$, and are orthogonal with respect to a
quasi-definite moment functional whose weights are modulated by the
Fox--Wright function $\Psi_\mu(z) = \sum_{k \ge 0}
\Gamma(k\mu{+}1)/\Gamma(k\mu{+}\mu{+}1)\cdot z^k$.

We place these polynomials precisely within the hierarchy of known orthogonal
families: they belong to the Nuttall--Stahl class of non-Hermitian orthogonal
polynomials \cite{NuttallSingh1977,Stahl1985I,Stahl1985II} on the Stahl
minimal-capacity compact $K_A(u)$, form a regular sequence in the sense of
Stahl--Totik \cite{StahlTotik1992}, and in the $\mu=1$ special case reduce
to Chebotarëv--Jacobi polynomials studied by
Barhoumi--Yattselev \cite{BarhoumiYattselev2023}. For $\mu \in (0,1)$ they
constitute a genuinely new subfamily -- \emph{Fox--Wright non-Hermitian
orthogonal polynomials} -- whose weight is a Fox--Wright-deformed
Jacobi-type function on a transcendental arc system, and whose asymptotic
recurrence coefficients are determined by the equilibrium measure of $K_A(u)$
via the Gonchar--Rakhmanov $S$-property \cite{GoncharRakhmanov1987}.
The Müntz--Szász density theorem \cite{MuntzSzasz1914,Moragues2018} provides
the underlying qualitative completeness backdrop, while the
Chebyshev--Wheeler algorithm \cite{Wheeler1974,Gragg1972} provides the
constructive computational realization.
\end{abstract}

\tableofcontents

% ---------------------------------------------------------------
\section{Introduction}
% ---------------------------------------------------------------

The Fractional Riccati--Müntz--Padé theorem \cite{Crowley2026} establishes
that the unique solution of the fractional Riccati equation
\begin{equation}\label{eq:fRic}
  D^{\mu} y(t,u) = P(u) + Q(u)\,y(t,u) + R(u)\,y(t,u)^{2},
  \qquad I^{1-\mu}y(0,u) = 0,
\end{equation}
with polynomial coefficients $P, Q, R \in \C[u]$, order $\mu \in (0,1]$,
and Riemann--Liouville calculus, is recovered in logarithmic capacity on
the complement of Stahl's compact $K_A(u)$ by the near-diagonal
Müntz--Padé approximants
\begin{equation}\label{eq:yM}
  y_M(t,u) = t^\mu \frac{N_M(t^\mu, u)}{D_M(t^\mu, u)}.
\end{equation}

Central to the proof is the identification of the Padé denominator sequence
$\{D_M\}$ with the orthogonal polynomial sequence of the moment functional
$\mathcal{L}[x^k] = a(k,u)$, where $\{a(k,u)\}$ are the Müntz--Tau
coefficients satisfying an explicit Fox--Wright-weighted recurrence
\cite{Crowley2026}. The present article characterizes this orthogonal
polynomial sequence in depth: its algebraic structure, its place in the
classification of orthogonal polynomial families, its asymptotic behavior,
and its relationship to classical and modern results in non-Hermitian
approximation theory.

The main contribution is the identification of these polynomials as
\emph{Fox--Wright non-Hermitian orthogonal polynomials} (FW-NHOPs) --
a new subfamily of the Nuttall--Stahl class parametrized by $\mu \in (0,1)$
that interpolates between the classical algebraic (Chebotarëv--Jacobi)
case at $\mu = 1$ and a transcendental Fox--Wright-weighted limit as
$\mu \to 0^+$. The Müntz--Szász condition $\sum 1/(k\mu) = \infty$ for all
$\mu \in (0,1]$ guarantees the underlying density \cite{Moragues2018},
while the Chebyshev--Wheeler algorithm \cite{Wheeler1974} provides
numerically stable computation of the recurrence coefficients.

% ---------------------------------------------------------------
\section{The Moment Functional and its Fox--Wright Structure}
% ---------------------------------------------------------------

\subsection{The Müntz--Tau Moments}

The fractional Riccati equation \eqref{eq:fRic}, recast as the Volterra
integral equation $y = I^\mu[P + Qy + Ry^2]$ via the Riemann--Liouville
composition formula \cite{SKM1993}, admits a unique formal power series
solution in the Müntz basis $\{t^{(k+1)\mu}\}_{k \ge 0}$
\cite{Crowley2026}:
\begin{equation}\label{eq:y_expansion}
  y(t,u) = \sum_{k \ge 0} a(k,u)\,t^{(k+1)\mu},
\end{equation}
with $a(k,\,\cdot\,) \in \C[u]$ satisfying the recurrence
\begin{align}
  a(0,u) &= P(u)\,\psi_{0}, \label{eq:a0}\\
  a(k,u) &= \psi_{k}\!\left(Q(u)\,a(k-1,u)
    + R(u)\sum_{j+\ell=k-1} a(j,u)\,a(\ell,u)\right), \quad k\ge 1,
    \label{eq:ak}
\end{align}
where the Fox--Wright weights are
\begin{equation}\label{eq:psi}
  \psi_k := \frac{\Gamma(k\mu+1)}{\Gamma(k\mu+\mu+1)}, \qquad k \ge 0.
\end{equation}

The sequence $\{\psi_k\}$ is precisely the sequence of eigenvalues of the
fractional integral $I^\mu$ on the Müntz monomial basis: $I^\mu
t^{(k+1)\mu} = \psi_{k+1}\,t^{(k+2)\mu}$ \cite{Crowley2026}. By Stirling's
formula, $\psi_k \sim (k\mu)^{-\mu}$ as $k\to\infty$, and $\Psi_\mu(z)
= \sum_{k\ge 0}\psi_k z^k$ has radius of convergence exactly $1$.

\subsection{Gevrey-1 Growth and Formal vs.\ Analytic Structure}

The generating function $A(z,u) = \sum_{k\ge 0} a(k,u)\,z^k$ lives in
$\C[u][[z]]$ as a formal power series. By an explicit induction on the
recurrence \eqref{eq:ak}, the moments satisfy the Gevrey-1 bound
\cite{Crowley2026}:
\begin{equation}\label{eq:gevrey}
  |a(k,u)| \le K_G(u)\,C_G(u)^k\,k!
\end{equation}
with computable constants $K_G, C_G$ depending on $|P(u)|$, $|Q(u)|$,
$|R(u)|$, and $\psi_0 = 1/\Gamma(\mu+1)$. The factorial growth means the
radius of convergence of $A$ is zero for generic parameters -- $A$ diverges
everywhere except at the origin.

This divergence is not a defect: it is the signature of the meromorphic
nature of $A$. Via a Mellin--Barnes integral representation \cite{Crowley2026}
\begin{equation}\label{eq:MB}
  A(z,u) = \frac{1}{2\pi i}\int_{c-i\infty}^{c+i\infty}
  M_{\Psi_\mu}(s)\,\widehat{B}_R(-s,u)\,(-z)^{-s}\,ds,
\end{equation}
where $M_{\Psi_\mu}(s) = \Gamma(s)\Gamma(1-s)\Gamma(1-s\mu)/\Gamma(\mu{+}1-s\mu)$
is the Mellin transform of $\Psi_\mu(-x)$, the formal series $A$ is promoted
to a meromorphic function on the universal cover $\widetilde{X}$ of $\C
\setminus \{0\}$ minus a minimal-capacity compact $K_A(u)$ -- Stahl's compact.

\subsection{The Hadamard Fixed-Point Equation}

Define $B[A](z,u) := P(u) + zQ(u)A(z,u) + zR(u)A(z,u)^2 \in \C[u][[z]]$.
Then the unique formal power series satisfying the Müntz--Tau recurrence is
characterized by the Hadamard fixed-point equation \cite{Crowley2026}:
\begin{equation}\label{eq:hadamard_fp}
  A(z,u) = \Psi_\mu(z) \odot B[A](z,u),
\end{equation}
where $\odot$ denotes the Hadamard (coefficient-wise) product:
$[z^k](F\odot G) = f_k g_k$. This identity holds in $\C[u][[z]]$ and is
proved by direct coefficient extraction \cite{Crowley2026}: at degree $k$,
$[z^k]B[A] = Q\,a(k{-}1) + R\sum_{j+\ell=k-1}a(j)a(\ell)$, and
multiplying by $\psi_k$ recovers $a(k)$ exactly from \eqref{eq:ak}.

% ---------------------------------------------------------------
\section{The Moment Functional and Its Quasi-Definiteness}
% ---------------------------------------------------------------

\begin{definition}[Fox--Wright moment functional]
The \emph{Fox--Wright moment functional} associated to \eqref{eq:fRic} is
the $\C$-linear map $\mathcal{L}: \C[x] \to \C[u]$ defined by
\begin{equation}\label{eq:L_def}
  \mathcal{L}[x^k] := a(k,u), \qquad k \ge 0,
\end{equation}
extended by linearity. The Hankel determinants are
\begin{equation}\label{eq:hankel}
  H_n(u) := \det\!\bigl(a(i+j,u)\bigr)_{0\le i,j\le n-1} \in \C[u].
\end{equation}
The functional is \emph{quasi-definite} if $H_n(u) \ne 0$ for all $n \ge 1$
(as elements of $\C[u]$, i.e., not identically zero).
\end{definition}

\begin{remark}[Generic quasi-definiteness]
By the Müntz--Tau recurrence \eqref{eq:a0}--\eqref{eq:ak}, each $a(k,u)$ is
a polynomial in $u$ of degree increasing with $k$. The Hankel determinant
$H_n(u)$ is a polynomial in $u$ that vanishes only on a proper algebraic
subvariety of parameter space; generically, $H_n(u) \ne 0$.
The recurrence coefficients $\alpha(k,u)$ and $\beta(k,u)$ produced by the
Chebyshev--Wheeler algorithm \cite{Wheeler1974,Gragg1972} are rational
functions of $u$ well-defined off this subvariety.
\end{remark}

Under quasi-definiteness, the Chebyshev--Wheeler algorithm
\cite{Wheeler1974,Gautschi1982} produces the unique monic polynomial
sequence $\{P_M(x,u)\}_{M\ge 0}$ satisfying
\begin{equation}\label{eq:orth_cond}
  \mathcal{L}[x^k P_M(x,u)] = 0, \qquad k = 0, 1, \dots, M-1,
\end{equation}
via the three-term recurrence
\begin{equation}\label{eq:three_term}
  P_{M+1}(x,u) = (x - \alpha(M,u))\,P_M(x,u) - \beta(M,u)\,P_{M-1}(x,u).
\end{equation}

The Padé denominator of $A(z,u)$ at order $M$ is the polynomial
\begin{equation}\label{eq:pade_denom}
  D_M(z,u) = z^M P_M(z^{-1},u),
\end{equation}
the reciprocal transform of $P_M$ \cite{Crowley2026,Gragg1972}.

% ---------------------------------------------------------------
\section{Classification: The Fox--Wright Non-Hermitian Orthogonal Polynomials}
% ---------------------------------------------------------------

\subsection{The Hierarchy of Orthogonal Polynomial Families}

We locate $\{P_M(\cdot, u)\}$ within the classification of orthogonal
polynomial sequences:

\begin{enumerate}
\item \textbf{Classical orthogonal polynomials} (Jacobi, Laguerre,
  Hermite) are orthogonal with respect to a positive measure on $\R$
  with explicit weight function. The moments grow at most polynomially.
  The recurrence coefficients $\alpha(k)$ and $\beta(k)$ are rational
  functions of $k$ with known explicit forms.
  \emph{Our polynomials do not belong to this class}: by \eqref{eq:gevrey},
  the moments grow like $k!$, ruling out any positive measure on $\R$
  \cite{StahlTotik1992}.

\item \textbf{Favard class / positive-measure OPs} are orthogonal with
  respect to a positive Borel measure on $\C$ with compact support.
  \emph{Our polynomials do not belong to this class} in general: the
  moment functional $\mathcal{L}$ is quasi-definite but not positive-definite
  (the Hankel determinants may have varying signs or be complex-valued for
  generic $u \in \C$).

\item \textbf{Regular OPs in the Stahl--Totik sense} \cite{StahlTotik1992}
  satisfy the $n$-th root asymptotics
  \begin{equation}\label{eq:nth_root}
    |P_M(z,u)|^{1/M} \;\to\; e^{g(z,K_A(u))}\,\mathrm{cap}(K_A(u))
    \quad \text{in capacity on } \C \setminus K_A(u),
  \end{equation}
  where $g(\cdot, K_A(u))$ is the Green function of $\C \setminus K_A(u)$
  with pole at infinity, and $\mathrm{cap}(K_A(u))$ is the logarithmic
  capacity. \emph{Our polynomials belong to this class} by virtue of
  Stahl's convergence theorem \cite{Stahl1997} as invoked in
  \cite{Crowley2026}.

\item \textbf{Nuttall--Stahl non-Hermitian OPs} \cite{NuttallSingh1977,
  Stahl1985I,Stahl1985II,GoncharRakhmanov1987} are orthogonal on the
  Stahl minimal-capacity contour $K_\Phi$ with a complex-valued weight
  $\rho$ (the jump of the approximated function across $K_\Phi$):
  \begin{equation}\label{eq:nh_orth}
    \int_{K_A(u)} z^k P_M(z,u)\,\rho(z,u)\,dz = 0,
    \qquad k = 0, \dots, M-1.
  \end{equation}
  \emph{Our polynomials belong to this class}, with $K_A(u)$ the Stahl
  compact from \cite{Crowley2026} and $\rho(z,u)$ the jump of $A(z,u)$
  across $K_A(u)$ -- a Fox--Wright-deformed Jacobi-type function described
  in Section~\ref{sec:weight} below.
\end{enumerate}

\subsection{The Weight Function}
\label{sec:weight}

The jump of $A(z,u)$ across $K_A(u)$ is determined by the algebraic
structure of the auxiliary function $F(z,u)$ satisfying the quadratic
\cite{Crowley2026}:
\begin{equation}\label{eq:quad}
  z R(u)\,F^2 - (1 - z Q(u))\,F + P(u)\,\Psi_\mu(z) = 0.
\end{equation}

The discriminant is
\begin{equation}\label{eq:disc}
  \Delta_F(z,u) = (1 - zQ(u))^2 - 4z P(u) R(u)\,\Psi_\mu(z),
\end{equation}
and $F(z,u)$ has the closed form
\begin{equation}\label{eq:F_closed}
  F(z,u) = \frac{(1-zQ) - \sqrt{\Delta_F}}{2zR}.
\end{equation}

The jump of $F$ across the branch cut $K_F(u)$ is
\begin{equation}\label{eq:jump_F}
  \rho_F(z,u) = F^+(z,u) - F^-(z,u) = \frac{\sqrt{\Delta_F^+(z,u)} - \sqrt{\Delta_F^-(z,u)}}{2zR(u)},
\end{equation}
which is a Fox--Wright-modified Jacobi-type weight: the classical Jacobi
weight $\sqrt{(z-e_1)(z-e_2)}$ (where $e_1, e_2$ are the branch points of
the algebraic part $(1-zQ)^2 - 4zPR$) is deformed by the factor $\Psi_\mu(z)$
in \eqref{eq:disc}. Since $A$ and $F$ have the same branch locus by the
branch-transfer lemma \cite{Crowley2026}, $\rho_A = \rho_F$ up to a
meromorphic factor, and the non-Hermitian orthogonality \eqref{eq:nh_orth}
holds with weight $\rho(z,u) = \rho_F(z,u)$.

\subsection{The Fox--Wright Deformation}

The key distinction from all previously named orthogonal polynomial families
is the presence of $\Psi_\mu(z)$ in the discriminant \eqref{eq:disc}. For
$\mu = 1$, one computes
\begin{equation}\label{eq:Psi1}
  \Psi_1(z) = \sum_{k\ge 0} \frac{1}{k+1}\,z^k = -\frac{\log(1-z)}{z},
\end{equation}
which is analytic in $|z| < 1$ but has a logarithmic branch point at $z=1$.
Thus even in the simplest non-classical case, $\Delta_F$ acquires a
transcendental factor. For $\mu \in (0,1)$, $\Psi_\mu$ is a Fox--Wright
${}_1\Psi_1$ function (see \cite{Wright1935,Gorenflo2014}):
\begin{equation}\label{eq:fox_wright_explicit}
  \Psi_\mu(z) = {}_1\Psi_1\!\left(\begin{matrix}(1,\mu)\\(1+\mu,\mu)\end{matrix}\,;\,z\right)
  = \sum_{k\ge 0}\frac{\Gamma(k\mu+1)}{\Gamma(k\mu+\mu+1)}\,z^k,
\end{equation}
with Mellin transform $M_{\Psi_\mu}(s) =
\Gamma(s)\Gamma(1-s)\Gamma(1-s\mu)/\Gamma(\mu{+}1-s\mu)$ meromorphic on $\C$
\cite{Crowley2026}. The branch points of $\sqrt{\Delta_F}$ are therefore the
zeros of \eqref{eq:disc}, which include both the zeros of the polynomial factor
$(1-zQ)^2 - 4zPR\cdot\psi_0$ (near $z=0$) and the singularities of
$\Psi_\mu$ on the circle $|z|=1$, projected under the Stahl--Chebotarëv
minimization \cite{Stahl1985I,Stahl1985II}.

% ---------------------------------------------------------------
\section{Asymptotic Theory}
% ---------------------------------------------------------------

\subsection{The $S$-Property and Stahl's Compact}

The Stahl compact $K_A(u)$ satisfies the \emph{$S$-property} of
Gonchar--Rakhmanov \cite{GoncharRakhmanov1987}: the normal derivative of
the Green function $g(z, K_A(u))$ from the two sides of $K_A(u)$ are equal
\begin{equation}\label{eq:S_property}
  \frac{\partial g}{\partial n^+}(z, K_A(u)) = \frac{\partial g}{\partial n^-}(z, K_A(u)),
  \qquad z \in K_A(u) \setminus \partial K_A(u).
\end{equation}
This characterizes $K_A(u)$ as the support of the equilibrium measure that
balances the logarithmic energy functional \cite{Ransford1995}:
\begin{equation}\label{eq:equil_meas}
  \nu_{K_A(u)} = \arg\min_\nu I[\nu], \qquad
  I[\nu] := \iint \log\frac{1}{|z-w|}\,d\nu(z)\,d\nu(w).
\end{equation}

\subsection{Zero Condensation: Nuttall--Stahl}

The empirical zero-counting measure of the orthogonal polynomials
\begin{equation}\label{eq:nu_M}
  \nu_M := \frac{1}{M}\sum_{P_M(\lambda,u)=0}\delta_\lambda
\end{equation}
converges weak-$*$ to $\nu_{K_A(u)}$ as $M\to\infty$, the equilibrium
measure of the Stahl compact \cite{Crowley2026,Stahl1997}. This is the
Nuttall--Stahl zero-condensation theorem, which in the present context
reads: the zeros of $\{P_M(\cdot, u)\}_{M\ge 1}$ are dense in $K_A(u)$
and equidistribute according to $\nu_{K_A(u)}$.

\subsection{Recurrence Coefficient Asymptotics}

Under the regularity hypothesis \eqref{eq:nth_root}, the recurrence
coefficients of \eqref{eq:three_term} satisfy \cite{StahlTotik1992}:
\begin{equation}\label{eq:alpha_asymp}
  \alpha(k, u) \;\longrightarrow\; \alpha^*(u), \qquad
  \beta(k, u) \;\longrightarrow\; \beta^*(u)
\end{equation}
as $k \to \infty$, where $\alpha^*(u)$ and $\beta^*(u)$ are determined by
the conformal map $\Phi: \C\setminus K_A(u) \to \C\setminus\overline{\mathbb{D}}$
by the Szegő-type asymptotic formula. Explicitly, the equilibrium measure
$\nu_{K_A(u)}$ has density $d\nu_{K_A(u)} = (1/2\pi)|\Phi'(z)|\,|dz|$ on
$K_A(u)$, and the asymptotic recurrence coefficients are the first two
Laurent coefficients of the conformal map:
\begin{equation}\label{eq:alpha_beta_conformal}
  \Phi(z) = z/\mathrm{cap}(K_A(u)) + \alpha^*(u) + O(1/z).
\end{equation}

For $K_A(u)$ an interval $[a,b] \subset \R$ (possible for special real
parameter values), this reduces to
\begin{equation}\label{eq:interval_case}
  \alpha^*(u) = \tfrac{a+b}{2}, \qquad \beta^*(u) = \tfrac{(b-a)^2}{16},
\end{equation}
and the $P_M(\cdot,u)$ are asymptotically Chebyshev polynomials for $[a,b]$.
For $K_A(u)$ a complex arc, the conformal map is transcendental and no
closed-form expression is known in general.

\subsection{Convergence: Stahl's Theorem Applied}

The Padé approximants $N_M/D_M$ converge to $A(z,u)$ in logarithmic
capacity on compact subsets of $\widetilde{X}\setminus K_A(u)$:
\begin{equation}\label{eq:stahl_conv}
  \left|\frac{N_M}{D_M} - A\right|^{1/(2M)} \;\longrightarrow\;
  e^{-g(z,K_A(u))} \quad \text{in capacity}
\end{equation}
\cite{Stahl1997,Crowley2026}. This implies exponential convergence off
$K_A(u)$ with rate governed by the Green function. Since $K_A(u)$ is
disjoint from $\R_+$ (the poles of $A$ on the physical time axis are
absent by a Volterra-iteration argument \cite{Crowley2026}), the convergence
upgrades to uniform convergence:
\begin{equation}\label{eq:uniform_conv}
  \sup_{t\in[0,T]}\bigl|y(t,u) - y_M(t,u)\bigr| \;\longrightarrow\; 0
  \qquad \text{as } M\to\infty.
\end{equation}

% ---------------------------------------------------------------
\section{The $\mu = 1$ Special Case: Chebotarëv--Jacobi Polynomials}
% ---------------------------------------------------------------

When $\mu = 1$, the Fox--Wright weights reduce to $\psi_k = 1/(k+1)$, and
$\Psi_1(z) = -\log(1-z)/z$. The discriminant \eqref{eq:disc} becomes
\begin{equation}\label{eq:disc_mu1}
  \Delta_F(z,u)\big|_{\mu=1} = (1-zQ)^2 + 4zPR\,\frac{\log(1-z)}{z}.
\end{equation}

The zeros of $\Delta_F$ are the solutions to $(1-zQ)^2 = -4PR\log(1-z)$,
a transcendental equation, but the branch structure of $\sqrt{\Delta_F}$ is
still two-sheeted, and the Stahl compact $K_A(u)$ is a finite arc in
$\C$ connecting the two branch points of $\sqrt{\Delta_F}$. The orthogonal
polynomials in this case are exactly the Padé denominator polynomials for a
two-sheeted function with one finite branch cut, the class studied in
\cite{BarhoumiYattselev2023} and earlier by Aptekarev--Stahl
\cite{Stahl1997}.

For the further special case $Q = 0$, $R = -1$, $P = 1$ (which gives the
Bernoulli equation at $\mu=1$), the quadratic \eqref{eq:quad} simplifies
and $K_A(u)$ is a segment $[0, \sigma]$ for $\sigma$ the first zero of the
classical Riccati solution -- the $P_M(\cdot,u)$ are then Legendre polynomials
shifted to $[0, \sigma]$.

% ---------------------------------------------------------------
\section{The $R = 0$ Oracle: Mittag-Leffler Moments and Laguerre Structure}
% ---------------------------------------------------------------

When $R(u) = 0$, the recurrence \eqref{eq:ak} linearizes to
$a(k,u) = \psi_k Q(u)^k P(u)\psi_0^{-1}\cdot\psi_0$, giving
\cite{Crowley2026}:
\begin{equation}\label{eq:r0_moments}
  a(k,u) = P(u)\,\psi_k\,Q(u)^k.
\end{equation}
The moment functional becomes $\mathcal{L}[x^k] = P(u)\psi_k Q(u)^k$, a
Stieltjes-type functional with moments $\{P\psi_k Q^k\}$.

\begin{proposition}
In the $R=0$ case, the moment functional $\mathcal{L}$ has the
integral representation
\begin{equation}\label{eq:r0_int}
  \mathcal{L}[f] = \int_0^\infty f(Qx)\,w_\mu(x)\,dx
\end{equation}
where $w_\mu(x) = P\,\mathcal{L}^{-1}[\psi_k](x)$ is the inverse Laplace
transform of the Fox--Wright function $z\mapsto P\,\Psi_\mu(Qz)$.
Explicitly, $w_\mu$ is a one-sided stable density with stability index $\mu$,
related to the Mittag-Leffler measure \cite{Grothaus2015}.
\end{proposition}

The $R=0$ orthogonal polynomials are therefore classical Laguerre-type
polynomials for the one-sided stable distribution at index $\mu$. For $\mu=1$,
$w_1(x) = P e^{-x}$ and the polynomials are standard Laguerre polynomials
$L_M^{(0)}$ scaled to $[0, 1/Q]$. For $\mu \ne 1$, they are the orthogonal
polynomials of the Mittag-Leffler measure of \cite{Grothaus2015} -- a family
studied in connection with fractional Brownian motion and anomalous diffusion.

This provides a closed-form unit test for any implementation: the $[M/M]$
Padé approximant to $A(z,u)|_{R=0} = P/(1-Qz)$ must reproduce the $M$-th
Padé approximant to $E_{\mu,1}(Qt^\mu)$ exactly, and the Wheeler algorithm
must recover Laguerre recurrence coefficients asymptotically \cite{Crowley2026}.

% ---------------------------------------------------------------
\section{The Müntz--Szász Backbone}
% ---------------------------------------------------------------

The Müntz--Szász theorem \cite{MuntzSzasz1914,Moragues2018} provides the
qualitative underpinning of the entire convergence theory. The Müntz
exponents used in the expansion \eqref{eq:y_expansion} are $\lambda_k =
(k+1)\mu$, and
\begin{equation}\label{eq:muntz_sum}
  \sum_{k=0}^\infty \frac{1}{\lambda_k} = \frac{1}{\mu}\sum_{k=0}^\infty \frac{1}{k+1} = \infty
\end{equation}
for every $\mu \in (0,1]$. By Müntz--Szász, the system $\{t^{(k+1)\mu}\}$
is therefore dense in $\{f \in C([0,T]) : f(0) = 0\}$ for all $T > 0$.

Three implications follow immediately:

\begin{enumerate}
\item \textbf{No coefficient can be discarded.} Any subsequence
  $\{a(k_j, u)\}$ with $\sum 1/k_j < \infty$ generates a Padé table that
  cannot uniformly recover $y(t,u)$ on $[0,T]$. All moments are
  indispensable.

\item \textbf{The Müntz rational approximants have maximum approximation power.}
  Since the exponents satisfy \eqref{eq:muntz_sum}, the Müntz rationals
  $\{t^\mu N_M(t^\mu, u)/D_M(t^\mu, u)\}$ form a complete approximating
  family in $C_0([0,T])$.

\item \textbf{Uniform convergence on $\R_+$ is the canonical upgrade.}
  The Stahl convergence-in-capacity result on the complement of $K_A(u)$,
  combined with the absence of poles on $\R_+$ (Lemma 14 of
  \cite{Crowley2026}), is the quantitative realization of the qualitative
  Müntz--Szász density -- the divergence of $\sum 1/\lambda_k$ is the
  engine, and the Fox--Wright-weighted Padé table is the mechanism.
\end{enumerate}

% ---------------------------------------------------------------
\section{Computational Realization: The Chebyshev--Wheeler Algorithm}
% ---------------------------------------------------------------

The three-term recurrence coefficients $\alpha(k,u)$ and $\beta(k,u)$ are
computed from the moments $\{a(k,u)\}$ by the Chebyshev--Wheeler algorithm
\cite{Wheeler1974,Gautschi1982}, which operates on an auxiliary array
$\sigma(k,\ell;u)$ initialized by $\sigma(0,\ell;u) = a(\ell,u)$ and
evolved by
\begin{equation}\label{eq:sigma_rec}
  \sigma(k,\ell;u) = \sigma(k-1,\ell+1;u) - \alpha(k-1,u)\,\sigma(k-1,\ell;u)
    - \beta(k-1,u)\,\sigma(k-2,\ell;u).
\end{equation}
The recurrence coefficients are
\begin{equation}\label{eq:alpha_beta_wheeler}
  \alpha(k,u) = \frac{\sigma(k,k+1;u)}{\sigma(k,k;u)}
    - \frac{\sigma(k-1,k;u)}{\sigma(k-1,k-1;u)},
  \qquad
  \beta(k,u) = \frac{\sigma(k,k;u)}{\sigma(k-1,k-1;u)}.
\end{equation}

Under quasi-definiteness, $\sigma(k,k;u) \ne 0$ for all $k$ and the
algorithm is well-defined. The orthonormality of the resulting polynomials
is guaranteed by Wheeler's theorem \cite{Wheeler1974}:
\begin{equation}\label{eq:wheeler_norm}
  \mathcal{L}[\pi_j\,\pi_k] = \delta_{jk}\prod_{i=0}^{k}\beta(i,u),
\end{equation}
where $\pi_k$ are the monic orthogonal polynomials produced by
\eqref{eq:three_term}.

\medskip
\noindent\textbf{Numerical stability.}
The Chebyshev--Wheeler algorithm for the Fox--Wright moment functional is
subject to the same conditioning issues as the general modified-moment
algorithm \cite{Gautschi1994,Gautschi1998}: the map from moments to
recurrence coefficients is a nonlinear map that can be ill-conditioned when
the moments grow rapidly. The Gevrey-1 bound \eqref{eq:gevrey} implies
moments of size $\sim K_G C_G^k k!$, and working in extended precision (as
in the ARB library \cite{Johansson2017}) is recommended for large $M$.

% ---------------------------------------------------------------
\section{Summary: Classification of the FW-NHOP Family}
% ---------------------------------------------------------------

\begin{table}[h]
\centering
\renewcommand{\arraystretch}{1.4}
\begin{tabular}{lll}
\hline
\textbf{Property} & \textbf{Value/Status} & \textbf{Reference} \\
\hline
Moment growth & Gevrey-1: $|a(k)| \le K C^k k!$ & \cite{Crowley2026} \\
Moment functional & Quasi-definite, not positive & \cite{Crowley2026} \\
Orthogonality & Non-Hermitian on $K_A(u)$ & \cite{NuttallSingh1977} \\
Weight & Fox--Wright-deformed Jacobi & \cite{Wright1935,Crowley2026} \\
Zero distribution & Condenses to $\nu_{K_A(u)}$ & \cite{Stahl1997,Crowley2026} \\
Regularity & Stahl--Totik regular & \cite{StahlTotik1992} \\
$S$-property & Gonchar--Rakhmanov & \cite{GoncharRakhmanov1987} \\
$\mu=1$ reduction & Chebotarëv--Jacobi & \cite{BarhoumiYattselev2023} \\
$R=0$ reduction & Mittag-Leffler / Laguerre-type & \cite{Grothaus2015} \\
Computation & Chebyshev--Wheeler & \cite{Wheeler1974,Gautschi1982} \\
Convergence & Stahl cap.\ + uniform on $\R_+$ & \cite{Stahl1997,Crowley2026} \\
Completeness & Müntz--Szász & \cite{Moragues2018} \\
\hline
\end{tabular}
\caption{Properties of the Fox--Wright non-Hermitian orthogonal polynomials
  (FW-NHOPs) associated to the fractional Riccati equation \eqref{eq:fRic}.}
\end{table}

The Fox--Wright non-Hermitian orthogonal polynomials form a one-parameter
family indexed by $\mu \in (0,1]$ that interpolates between:
\begin{itemize}
  \item $\mu = 1$: classical (logarithmically-deformed) Chebotarëv--Jacobi
    polynomials on a two-arc system \cite{BarhoumiYattselev2023};
  \item $\mu \to 0^+$: a limiting regime where $\psi_k \to 1$ for all $k$
    (by $\Gamma(1)/\Gamma(1) = 1$) and the Fox--Wright weight collapses to
    a standard algebraic weight -- the polynomials approach classical
    Padé denominators for the bare quadratic $F_0(z,u)$ with $\Psi_0 \equiv 1$.
\end{itemize}

% ---------------------------------------------------------------
\begin{thebibliography}{99}

\bibitem{Crowley2026}
S.~Crowley,
\emph{The Fractional Riccati--Müntz--Padé Theorem},
preprint (2026).

\bibitem{SKM1993}
S.~G.~Samko, A.~A.~Kilbas, O.~I.~Marichev,
\emph{Fractional Integrals and Derivatives: Theory and Applications},
Gordon and Breach, 1993.

\bibitem{Stahl1985I}
H.~Stahl,
\emph{Extremal domains associated with an analytic function I},
Constr.\ Approx.\ \textbf{1} (1985), 89--110.

\bibitem{Stahl1985II}
H.~Stahl,
\emph{Extremal domains associated with an analytic function II},
Constr.\ Approx.\ \textbf{1} (1985), 111--137.

\bibitem{Stahl1997}
H.~Stahl,
\emph{The convergence of Padé approximants to functions with branch points},
J.\ Approx.\ Theory \textbf{91} (1997), 139--204.

\bibitem{StahlTotik1992}
H.~Stahl, V.~Totik,
\emph{General Orthogonal Polynomials},
Encyclopedia of Mathematics and its Applications \textbf{43},
Cambridge University Press, 1992.

\bibitem{GoncharRakhmanov1987}
A.~A.~Gonchar, E.~A.~Rakhmanov,
\emph{Equilibrium distributions and degree of rational approximation of
analytic functions},
Mat.\ Sb.\ \textbf{134}(176) (1987), 306--352;
English transl.\ Math.\ USSR-Sb.\ \textbf{62} (1989), 305--348.

\bibitem{NuttallSingh1977}
J.~Nuttall, S.~R.~Singh,
\emph{Orthogonal polynomials and Padé approximants associated with a system
of arcs},
J.\ Approx.\ Theory \textbf{21} (1977), 1--42.

\bibitem{BarhoumiYattselev2023}
A.~B.~Barhoumi, M.~L.~Yattselev,
\emph{Non-Hermitian orthogonal polynomials on a trefoil},
arXiv:2303.17037 (2023).

\bibitem{Wheeler1974}
J.~C.~Wheeler,
\emph{Modified moments and Gaussian quadratures},
Rocky Mountain J.\ Math.\ \textbf{4} (1974), 287--296.

\bibitem{Gragg1972}
W.~B.~Gragg,
\emph{The Padé table and its relation to certain algorithms of numerical
analysis},
SIAM Review \textbf{14} (1972), 1--62.

\bibitem{Gautschi1982}
W.~Gautschi,
\emph{On generating orthogonal polynomials},
SIAM J.\ Sci.\ Stat.\ Comput.\ \textbf{3} (1982), 289--317.

\bibitem{Gautschi1994}
W.~Gautschi,
\emph{Orthogonal polynomials: applications and computation},
Acta Numerica \textbf{5} (1996), 45--119.

\bibitem{Gautschi1998}
W.~Gautschi, L.~Gori, M.~L.~Lo Cascio,
\emph{Quadrature rules for rational functions},
Numer.\ Math.\ \textbf{86} (2000), 617--633.

\bibitem{MuntzSzasz1914}
Ch.~H.~Müntz,
\emph{Über den Approximationssatz von Weierstrass},
H.~A.~Schwarz Festschrift, Berlin, 1914, 303--312.

\bibitem{Moragues2018}
A.~Ferré~Moragues,
\emph{What is\ldots\ the Müntz--Szász Theorem?},
expository note (2018).

\bibitem{Wright1935}
E.~M.~Wright,
\emph{The asymptotic expansion of the generalized Bessel function},
Proc.\ London Math.\ Soc.\ (2) \textbf{38} (1935), 257--270.

\bibitem{Gorenflo2014}
R.~Gorenflo, A.~A.~Kilbas, F.~Mainardi, S.~V.~Rogosin,
\emph{Mittag-Leffler Functions, Related Topics and Applications},
Springer, 2014.

\bibitem{Grothaus2015}
M.~Grothaus, F.~Jahnert,
\emph{Mittag-Leffler analysis I: Construction and characterization},
J.\ Funct.\ Anal.\ \textbf{268} (2015), 1876--1903.

\bibitem{Ransford1995}
T.~Ransford,
\emph{Potential Theory in the Complex Plane},
London Math.\ Soc.\ Student Texts \textbf{28}, Cambridge Univ.\ Press, 1995.

\bibitem{Johansson2017}
F.~Johansson,
\emph{Arb: efficient arbitrary-precision midpoint-radius interval arithmetic},
IEEE Trans.\ Comput.\ \textbf{66} (2017), 1281--1292.

\end{thebibliography}

\end{document}
